% Options for packages loaded elsewhere
\PassOptionsToPackage{unicode}{hyperref}
\PassOptionsToPackage{hyphens}{url}
\PassOptionsToPackage{dvipsnames,svgnames,x11names}{xcolor}
\documentclass[
  brazilian,
  11pt,
]{article}
\usepackage{xcolor}
\usepackage[margin=2.4cm]{geometry}
\usepackage{amsmath,amssymb}
\setcounter{secnumdepth}{-\maxdimen} % remove section numbering
\usepackage{iftex}
\ifPDFTeX
  \usepackage[T1]{fontenc}
  \usepackage[utf8]{inputenc}
  \usepackage{textcomp} % provide euro and other symbols
\else % if luatex or xetex
  \usepackage{unicode-math} % this also loads fontspec
  \defaultfontfeatures{Scale=MatchLowercase}
  \defaultfontfeatures[\rmfamily]{Ligatures=TeX,Scale=1}
\fi
\usepackage{lmodern}
\ifPDFTeX\else
  % xetex/luatex font selection
\fi
% Use upquote if available, for straight quotes in verbatim environments
\IfFileExists{upquote.sty}{\usepackage{upquote}}{}
\IfFileExists{microtype.sty}{% use microtype if available
  \usepackage[]{microtype}
  \UseMicrotypeSet[protrusion]{basicmath} % disable protrusion for tt fonts
}{}
\makeatletter
\@ifundefined{KOMAClassName}{% if non-KOMA class
  \IfFileExists{parskip.sty}{%
    \usepackage{parskip}
  }{% else
    \setlength{\parindent}{0pt}
    \setlength{\parskip}{6pt plus 2pt minus 1pt}}
}{% if KOMA class
  \KOMAoptions{parskip=half}}
\makeatother
\usepackage{color}
\usepackage{fancyvrb}
\newcommand{\VerbBar}{|}
\newcommand{\VERB}{\Verb[commandchars=\\\{\}]}
\DefineVerbatimEnvironment{Highlighting}{Verbatim}{commandchars=\\\{\}}
% Add ',fontsize=\small' for more characters per line
\usepackage{framed}
\definecolor{shadecolor}{RGB}{35,38,41}
\newenvironment{Shaded}{\begin{snugshade}}{\end{snugshade}}
\newcommand{\AlertTok}[1]{\textcolor[rgb]{0.58,0.85,0.30}{\textbf{\colorbox[rgb]{0.30,0.12,0.14}{#1}}}}
\newcommand{\AnnotationTok}[1]{\textcolor[rgb]{0.25,0.50,0.35}{#1}}
\newcommand{\AttributeTok}[1]{\textcolor[rgb]{0.16,0.50,0.73}{#1}}
\newcommand{\BaseNTok}[1]{\textcolor[rgb]{0.96,0.45,0.00}{#1}}
\newcommand{\BuiltInTok}[1]{\textcolor[rgb]{0.50,0.55,0.55}{#1}}
\newcommand{\CharTok}[1]{\textcolor[rgb]{0.24,0.68,0.91}{#1}}
\newcommand{\CommentTok}[1]{\textcolor[rgb]{0.48,0.49,0.49}{#1}}
\newcommand{\CommentVarTok}[1]{\textcolor[rgb]{0.50,0.55,0.55}{#1}}
\newcommand{\ConstantTok}[1]{\textcolor[rgb]{0.15,0.68,0.68}{\textbf{#1}}}
\newcommand{\ControlFlowTok}[1]{\textcolor[rgb]{0.99,0.74,0.29}{\textbf{#1}}}
\newcommand{\DataTypeTok}[1]{\textcolor[rgb]{0.16,0.50,0.73}{#1}}
\newcommand{\DecValTok}[1]{\textcolor[rgb]{0.96,0.45,0.00}{#1}}
\newcommand{\DocumentationTok}[1]{\textcolor[rgb]{0.64,0.20,0.25}{#1}}
\newcommand{\ErrorTok}[1]{\textcolor[rgb]{0.85,0.27,0.33}{\underline{#1}}}
\newcommand{\ExtensionTok}[1]{\textcolor[rgb]{0.00,0.60,1.00}{\textbf{#1}}}
\newcommand{\FloatTok}[1]{\textcolor[rgb]{0.96,0.45,0.00}{#1}}
\newcommand{\FunctionTok}[1]{\textcolor[rgb]{0.56,0.27,0.68}{#1}}
\newcommand{\ImportTok}[1]{\textcolor[rgb]{0.15,0.68,0.38}{#1}}
\newcommand{\InformationTok}[1]{\textcolor[rgb]{0.77,0.36,0.00}{#1}}
\newcommand{\KeywordTok}[1]{\textcolor[rgb]{0.81,0.81,0.76}{\textbf{#1}}}
\newcommand{\NormalTok}[1]{\textcolor[rgb]{0.81,0.81,0.76}{#1}}
\newcommand{\OperatorTok}[1]{\textcolor[rgb]{0.81,0.81,0.76}{#1}}
\newcommand{\OtherTok}[1]{\textcolor[rgb]{0.15,0.68,0.38}{#1}}
\newcommand{\PreprocessorTok}[1]{\textcolor[rgb]{0.15,0.68,0.38}{#1}}
\newcommand{\RegionMarkerTok}[1]{\textcolor[rgb]{0.16,0.50,0.73}{\colorbox[rgb]{0.08,0.19,0.26}{#1}}}
\newcommand{\SpecialCharTok}[1]{\textcolor[rgb]{0.24,0.68,0.91}{#1}}
\newcommand{\SpecialStringTok}[1]{\textcolor[rgb]{0.85,0.27,0.33}{#1}}
\newcommand{\StringTok}[1]{\textcolor[rgb]{0.96,0.31,0.31}{#1}}
\newcommand{\VariableTok}[1]{\textcolor[rgb]{0.15,0.68,0.68}{#1}}
\newcommand{\VerbatimStringTok}[1]{\textcolor[rgb]{0.85,0.27,0.33}{#1}}
\newcommand{\WarningTok}[1]{\textcolor[rgb]{0.85,0.27,0.33}{#1}}
\ifLuaTeX
\usepackage[bidi=basic,shorthands=off]{babel}
\else
\usepackage[bidi=default,shorthands=off]{babel}
\fi
\ifLuaTeX
  \usepackage{selnolig} % disable illegal ligatures
\fi
\setlength{\emergencystretch}{3em} % prevent overfull lines
\providecommand{\tightlist}{%
  \setlength{\itemsep}{0pt}\setlength{\parskip}{0pt}}
% ==========================================================================
%  Cabecalho LaTeX do curso "Trabalhando com Texto em Python I"
%  Injetado pelo pandoc (-H) ao gerar o .tex de cada unidade.
% ==========================================================================
\usepackage[T1]{fontenc}
\usepackage{lmodern}
\usepackage[scaled=0.92]{helvet}
\renewcommand{\familydefault}{\sfdefault}
\usepackage{microtype}
\usepackage{amssymb}
\usepackage{newunicodechar}
\usepackage{xcolor}
\usepackage[most]{tcolorbox}
\tcbuselibrary{skins,breakable}
\usepackage{titlesec}
\usepackage{enumitem}
% quebra linhas longas nos blocos de codigo (evita vazar a margem)
\usepackage{fvextra}
\fvset{breaklines=true,breakanywhere=true,fontsize=\small}

% -------- paleta --------
\definecolor{accent}{HTML}{6C4CF1}
\definecolor{cyan}{HTML}{00B4D8}
\definecolor{subink}{HTML}{2B3242}
\definecolor{notec}{HTML}{3B82F6}
\definecolor{tipc}{HTML}{10B981}
\definecolor{warnc}{HTML}{F59E0B}

% -------- emojis / simbolos unicode que podem aparecer no texto --------
% simbolos com significado
\newunicodechar{✅}{{\color{tipc}\checkmark}}
\newunicodechar{✔}{{\color{tipc}\checkmark}}
\newunicodechar{☐}{$\square$}
\newunicodechar{κ}{\ensuremath{\kappa}}
\newunicodechar{≈}{\ensuremath{\approx}}
\newunicodechar{−}{\ensuremath{-}}
\newunicodechar{↓}{\ensuremath{\downarrow}}
\newunicodechar{→}{\ensuremath{\rightarrow}}
\newunicodechar{←}{\ensuremath{\leftarrow}}
\newunicodechar{↹}{\ensuremath{\hookrightarrow}}
\newunicodechar{°}{\textdegree}
\newunicodechar{·}{\textperiodcentered}
\newunicodechar{´}{'}
\newunicodechar{‑}{-}
% emojis decorativos -> removidos no PDF
\newunicodechar{💡}{}
\newunicodechar{🚀}{}
\newunicodechar{📝}{}
\newunicodechar{⚠}{}
\newunicodechar{🐍}{}
\newunicodechar{🎉}{}
\newunicodechar{🍕}{}
\newunicodechar{👍}{}
\newunicodechar{💋}{}
\newunicodechar{😚}{}
\newunicodechar{🚂}{}
\newunicodechar{❤}{}

% -------- caixas de callout (usadas pelo filtro callouts.lua) --------
\newtcolorbox{notebox}{enhanced,breakable,colback=notec!6,colframe=notec,
  boxrule=0pt,leftrule=4pt,arc=3pt,left=10pt,right=10pt,top=7pt,bottom=7pt,
  before skip=10pt,after skip=12pt}
\newtcolorbox{tipbox}{enhanced,breakable,colback=tipc!7,colframe=tipc,
  boxrule=0pt,leftrule=4pt,arc=3pt,left=10pt,right=10pt,top=7pt,bottom=7pt,
  before skip=10pt,after skip=12pt}
\newtcolorbox{warnbox}{enhanced,breakable,colback=warnc!8,colframe=warnc,
  boxrule=0pt,leftrule=4pt,arc=3pt,left=10pt,right=10pt,top=7pt,bottom=7pt,
  before skip=10pt,after skip=12pt}
\newtcolorbox{objbox}{enhanced,breakable,colback=accent!5,colframe=accent,
  boxrule=0pt,leftrule=4pt,arc=3pt,left=10pt,right=10pt,top=7pt,bottom=8pt,
  before skip=10pt,after skip=14pt}

% -------- titulos coloridos --------
\titleformat{\section}{\sffamily\Large\bfseries\color{accent}}{\thesection}{0.6em}{}
\titleformat{\subsection}{\sffamily\large\bfseries\color{subink}}{\thesubsection}{0.5em}{}
\titleformat{\subsubsection}{\sffamily\normalsize\bfseries\color{cyan}}{\thesubsubsection}{0.5em}{}
\titlespacing*{\section}{0pt}{2.2ex plus 1ex minus .2ex}{1.1ex}

% codigo inline colorido
\definecolor{inlink}{HTML}{5A34D6}
\let\oldtexttt\texttt
\renewcommand{\texttt}[1]{{\color{inlink}\oldtexttt{#1}}}
\usepackage{bookmark}
\IfFileExists{xurl.sty}{\usepackage{xurl}}{} % add URL line breaks if available
\urlstyle{same}
% fallback for those not using the hyperref driver hyperxmp:
\makeatletter
\@ifundefined{xmpquote}{\newcommand{\xmpquote}[1]{#1}}{}
\makeatother
\hypersetup{
  pdftitle={Manipulando texto com Python},
  pdfauthor={CiberExt 26-29 · FEELT38103 · Universidade Federal de Uberlândia},
  pdflang={pt-BR},
  colorlinks=true,
  linkcolor={accent},
  filecolor={Maroon},
  citecolor={Blue},
  urlcolor={cyan},
  pdfcreator={LaTeX via pandoc}}

\title{Manipulando texto com Python}
\usepackage{etoolbox}
\makeatletter
\providecommand{\subtitle}[1]{% add subtitle to \maketitle
  \apptocmd{\@title}{\par {\large #1 \par}}{}{}
}
\makeatother
\subtitle{Trabalhando com Texto em Python I · Unidade 1}
\author{CiberExt 26-29 · FEELT38103 · Universidade Federal de
Uberlândia}
\date{2026}

\begin{document}
\maketitle

{
\setcounter{tocdepth}{2}
\tableofcontents
}
\subsubsection{Objetivos de
aprendizagem}\label{objetivos-de-aprendizagem}

Ao final desta unidade, você deverá ser capaz de:

\begin{itemize}
\tightlist
\item
  \textbf{compreender a diferença} entre texto formatado (\emph{rich
  text}), texto estruturado (\emph{structured text}) e texto simples
  (\emph{plain text});
\item
  \textbf{entender o conceito} de codificação de caracteres (\emph{text
  encoding});
\item
  \textbf{carregar arquivos} de texto simples no Python e
  \textbf{manipular} o seu conteúdo;
\item
  \textbf{definir padrões flexíveis} de busca com expressões regulares;
\item
  \textbf{processar vários arquivos} de uma só vez e salvar os
  resultados.
\end{itemize}

\begin{center}\rule{0.5\linewidth}{0.5pt}\end{center}

\subsection{1. Computadores e texto}\label{computadores-e-texto}

Os computadores podem armazenar e representar texto em formatos
diferentes. Conhecer a distinção entre esses tipos é
\textbf{fundamental} para processá‑los programaticamente, porque cada
formato traz consigo informações (ou ruídos) distintos.

\subsubsection{\texorpdfstring{1.1 O que é texto formatado (\emph{rich
text})?}{1.1 O que é texto formatado (rich text)?}}\label{o-que-uxe9-texto-formatado-rich-text}

Processadores de texto, como o Microsoft Word, produzem \textbf{texto
formatado} (\emph{rich text}): texto cuja aparência foi estilizada de
alguma maneira específica.

O texto formatado permite definir estilos visuais para os elementos do
documento. Títulos, por exemplo, podem usar uma fonte diferente da do
corpo do texto, que por sua vez pode empregar itálico ou negrito para
dar ênfase. O texto formatado também pode incluir imagens, tabelas e
outros elementos.

Esse é o formato padrão dos processadores de texto modernos do tipo
\emph{WYSIWYG} (\emph{what you see is what you get} --- ``o que você vê
é o que você obtém'').

\subsubsection{\texorpdfstring{1.2 O que é texto simples (\emph{plain
text})?}{1.2 O que é texto simples (plain text)?}}\label{o-que-uxe9-texto-simples-plain-text}

Diferentemente do texto formatado, o \textbf{texto simples} (\emph{plain
text}) não contém nenhuma informação sobre a aparência visual: ele é
feito \textbf{apenas de caracteres}.

Neste contexto, ``caracteres'' abrange letras, números, sinais de
pontuação, espaços e quebras de linha. A definição de texto simples é um
tanto flexível, mas, em geral, refere‑se a texto sem qualquer informação
de formatação ou estilo.

\subsubsection{\texorpdfstring{1.3 O que é texto estruturado
(\emph{structured
text})?}{1.3 O que é texto estruturado (structured text)?}}\label{o-que-uxe9-texto-estruturado-structured-text}

O \textbf{texto estruturado} (\emph{structured text}) pode ser entendido
como um caso especial de texto simples que inclui sequências de
caracteres usadas para \textbf{formatar o texto para exibição}.

São exemplos de texto estruturado os textos descritos por linguagens de
marcação como XML, Markdown ou HTML.

O exemplo abaixo mostra uma frase em texto simples envolvida por
marcadores (\emph{tags}) HTML de parágrafo
\texttt{\textless{}p\textgreater{}}. A marca de abertura
\texttt{\textless{}p\textgreater{}} e a de fechamento
\texttt{\textless{}/p\textgreater{}} informam ao computador que todo o
conteúdo colocado entre elas forma um parágrafo:

\begin{Shaded}
\begin{Highlighting}[]
\DataTypeTok{\textless{}}\KeywordTok{p}\DataTypeTok{\textgreater{}}\NormalTok{This is an example sentence.}\DataTypeTok{\textless{}/}\KeywordTok{p}\DataTypeTok{\textgreater{}}
\end{Highlighting}
\end{Shaded}

Essa informação é usada para estruturar o texto simples no momento de
renderizá‑lo para exibição, normalmente aplicando estilo à sua
aparência.

\subsubsection{1.4 Por que isso importa?}\label{por-que-isso-importa}

Ao reunir um conjunto de textos para formar um \textbf{corpus}, é bem
provável que parte deles tenha origem em formato formatado ou
estruturado, dependendo do meio de onde vieram:

\begin{itemize}
\tightlist
\item
  Se você coleta documentos impressos que foram digitalizados por
  \textbf{reconhecimento óptico de caracteres} (\emph{optical character
  recognition}, OCR) e depois convertidos de texto formatado para texto
  simples, a remoção das informações de formatação tende a
  \textbf{introduzir erros} no texto resultante. Trabalhar com esse tipo
  de OCR ``sujo'' pode afetar os resultados da análise textual (Hill \&
  Hengchen, 2019).
\item
  Se você coleta documentos digitais raspando fóruns de discussão ou
  sites, é provável que encontre vestígios de texto estruturado na forma
  de marcadores de marcação, que podem ser arrastados para o texto
  simples durante a conversão.
\end{itemize}

O texto simples é, de longe, o formato \textbf{mais interoperável} para
texto, por ser fácil de ler para os computadores. É por isso que as
linguagens de programação trabalham com texto simples --- e, se você
pretende usar programação para manipular texto, precisa saber o que é
texto simples.

\begin{notebox}

\textbf{Em resumo:} ao trabalhar com texto simples, muitas vezes será
preciso lidar com os vestígios deixados pela conversão a partir de texto
formatado ou estruturado.

\end{notebox}

\begin{center}\rule{0.5\linewidth}{0.5pt}\end{center}

\subsection{\texorpdfstring{2. Codificação de texto (\emph{text
encoding})}{2. Codificação de texto (text encoding)}}\label{codificauxe7uxe3o-de-texto-text-encoding}

Para serem lidos pelos computadores, os textos simples precisam ser
\textbf{codificados}. Isso é feito por meio da \textbf{codificação de
caracteres} (\emph{character encoding}), que mapeia cada caractere
(letras, números, pontuação, espaços\ldots) para uma representação
numérica compreendida pela máquina.

O ideal seria não termos de lidar com operações de baixo nível como a
codificação de caracteres --- mas, na prática, precisamos, porque
existem \textbf{vários sistemas} de codificação e eles \textbf{não são
compatíveis entre si}. Essa é a origem de boa parte das dores de cabeça
de quem trabalha com texto simples.

Há dois sistemas de codificação com os quais você provavelmente vai se
deparar: \textbf{ASCII} e \textbf{Unicode}.

\subsubsection{2.1 ASCII}\label{ascii}

O \textbf{ASCII} (\emph{American Standard Code for Information
Interchange}) é um sistema pioneiro de codificação de caracteres que
serviu de base para muitos sistemas modernos.

O ASCII ainda é bastante usado, mas é \textbf{muito limitado} quanto à
variedade de caracteres. Se o seu idioma inclui caracteres como
\texttt{ä}, \texttt{ö} --- ou os nossos \texttt{á}, \texttt{ç},
\texttt{ã} ---, o ASCII não dá conta.

\subsubsection{2.2 Unicode}\label{unicode}

O \textbf{Unicode} é um padrão para codificar texto na maioria dos
sistemas de escrita usados no mundo, cobrindo cerca de \textbf{140 mil
caracteres} entre escritas modernas e históricas, símbolos e emojis.

Por exemplo, o emoji de fatia de pizza 🍕 tem o ``código'' Unicode
\texttt{U+1F355}, enquanto o código correspondente a um espaço em branco
é \texttt{U+0020}.

O Unicode pode ser implementado por diferentes codificações, como o
\textbf{UTF‑8}, definido pelo próprio padrão Unicode.

O UTF‑8 é \textbf{retrocompatível com o ASCII}. Em outras palavras, as
codificações do ASCII formam um subconjunto do UTF‑8 --- o que facilita
bastante a nossa vida. Assim, mesmo que um arquivo de texto simples
tenha sido codificado em ASCII, podemos decodificá‑lo como UTF‑8;
\textbf{o contrário, porém, não vale}.

\begin{center}\rule{0.5\linewidth}{0.5pt}\end{center}

\subsection{3. Carregando arquivos de texto simples no
Python}\label{carregando-arquivos-de-texto-simples-no-python}

Arquivos de texto simples podem ser carregados no Python com a função
\texttt{open()}.

O primeiro argumento de \texttt{open()} deve ser uma \emph{string}
contendo o \textbf{caminho} para o arquivo que será aberto. Aqui, temos
um caminho que aponta para um arquivo chamado
\texttt{NYT\_1991-01-16-A15.txt}, localizado em um diretório chamado
\texttt{data}. Na definição do caminho, o diretório e o nome do arquivo
são separados por uma barra \texttt{/}.

Para acessar o arquivo apontado por esse caminho, passamos o caminho
como \emph{string} ao argumento \texttt{file} da função \texttt{open()}:

\begin{Shaded}
\begin{Highlighting}[]
\BuiltInTok{open}\NormalTok{(}\BuiltInTok{file}\OperatorTok{=}\StringTok{\textquotesingle{}data/NYT\_1991{-}01{-}16{-}A15.txt\textquotesingle{}}\NormalTok{, mode}\OperatorTok{=}\StringTok{\textquotesingle{}r\textquotesingle{}}\NormalTok{, encoding}\OperatorTok{=}\StringTok{\textquotesingle{}utf{-}8\textquotesingle{}}\NormalTok{)}
\end{Highlighting}
\end{Shaded}

Vale entender os demais argumentos:

\begin{itemize}
\tightlist
\item
  Por padrão, o Python 3 assume que o texto está codificado em
  \textbf{UTF‑8}, mas podemos deixar isso explícito com o argumento
  \texttt{encoding}, passando a \emph{string} \texttt{utf-8}.
\item
  O argumento \texttt{mode} define o que queremos fazer com o arquivo. A
  \emph{string} \texttt{r} (de \emph{read}) indica que queremos
  \textbf{apenas ler} o arquivo.
\item
  Usamos \texttt{open()} junto da instrução \texttt{with}, que garante
  que o arquivo será \textbf{fechado} após executarmos o que precisamos
  dentro do bloco de código indentado que segue o \texttt{with}. Isso
  evita que o arquivo continue consumindo memória e recursos depois que
  não precisamos mais dele.
\end{itemize}

\begin{Shaded}
\begin{Highlighting}[]
\CommentTok{\# Abre um arquivo e o atribui à variável \textquotesingle{}file\textquotesingle{}}
\ControlFlowTok{with} \BuiltInTok{open}\NormalTok{(}\BuiltInTok{file}\OperatorTok{=}\StringTok{\textquotesingle{}data/NYT\_1991{-}01{-}16{-}A15.txt\textquotesingle{}}\NormalTok{, mode}\OperatorTok{=}\StringTok{\textquotesingle{}r\textquotesingle{}}\NormalTok{, encoding}\OperatorTok{=}\StringTok{\textquotesingle{}utf{-}8\textquotesingle{}}\NormalTok{) }\ImportTok{as} \BuiltInTok{file}\NormalTok{:}

    \CommentTok{\# A instrução \textquotesingle{}with\textquotesingle{} deve ser seguida por um bloco de código indentado.}
    \CommentTok{\# Aqui chamamos o método read() para ler o conteúdo do arquivo e}
    \CommentTok{\# atribuímos o resultado à variável \textquotesingle{}text\textquotesingle{}.}
\NormalTok{    text }\OperatorTok{=} \BuiltInTok{file}\NormalTok{.read()}
\end{Highlighting}
\end{Shaded}

Repare que \texttt{with} e \texttt{open()} são seguidos da palavra
\texttt{as} e de uma variável chamada \texttt{file}. Isso instrui o
Python a atribuir à variável \texttt{file} aquilo que \texttt{open()}
retornar.

Se agora chamarmos a variável \texttt{file}, obtemos um objeto Python do
tipo \texttt{TextIOWrapper}, que carrega três informações: o caminho do
arquivo (no argumento \texttt{name}) e os argumentos \texttt{mode} e
\texttt{encoding} que definimos acima:

\begin{Shaded}
\begin{Highlighting}[]
\CommentTok{\# Chama a variável para examinar o objeto}
\BuiltInTok{file}
\end{Highlighting}
\end{Shaded}

\begin{Shaded}
\begin{Highlighting}[]
\NormalTok{\textless{}\_io.TextIOWrapper name=\textquotesingle{}data/NYT\_1991{-}01{-}16{-}A15.txt\textquotesingle{} mode=\textquotesingle{}r\textquotesingle{} encoding=\textquotesingle{}utf{-}8\textquotesingle{}\textgreater{}}
\end{Highlighting}
\end{Shaded}

Tenha em mente que, no bloco indentado após o \texttt{with}, chamamos o
método \texttt{read()} do objeto \texttt{TextIOWrapper}. Esse método leu
o conteúdo do arquivo, que atribuímos à variável \texttt{text}.

No entanto, se tentarmos chamar \texttt{read()} para a variável
\texttt{file} \textbf{fora} do bloco \texttt{with}, o Python levanta um
erro, porque o arquivo já foi fechado:

\begin{Shaded}
\begin{Highlighting}[]
\CommentTok{\# Tenta usar o método read() para ler o conteúdo do arquivo}
\BuiltInTok{file}\NormalTok{.read()}
\end{Highlighting}
\end{Shaded}

\begin{Shaded}
\begin{Highlighting}[]
\NormalTok{ValueError: I/O operation on closed file.}
\end{Highlighting}
\end{Shaded}

Esse comportamento é o esperado: queremos que o arquivo seja fechado,
para que não consuma memória nem recursos quando não é mais necessário.
Isso é especialmente importante ao trabalhar com \textbf{milhares de
arquivos}, pois cada arquivo aberto ocupa memória.

Vejamos agora o resultado de \texttt{read()}, guardado em \texttt{text}.
O texto é bem longo, então vamos pegar apenas uma \textbf{fatia} com os
primeiros 500 caracteres, usando colchetes \texttt{{[}:500{]}}.
Adicionar colchetes logo após o nome de uma variável permite acessar
partes do objeto, quando o objeto permite isso. Por exemplo,
\texttt{text{[}1{]}} recuperaria o caractere na posição 1; o \texttt{:}
antes do número instrui o Python a recuperar \textbf{todos} os
caracteres até o de posição 500.

\begin{Shaded}
\begin{Highlighting}[]
\CommentTok{\# Recupera os primeiros 500 caracteres da variável \textquotesingle{}text\textquotesingle{}}
\NormalTok{text[:}\DecValTok{500}\NormalTok{]}
\end{Highlighting}
\end{Shaded}

\begin{Shaded}
\begin{Highlighting}[]
\NormalTok{\textquotesingle{}U.S. TAKING STEPS TO CURB TERRORISM: F.B.I. Is Ordered to Find Iraqis Whose Visas Have Expired\textbackslash{}nBy JAMES BARRON\textbackslash{}nNew York Times (1923{-}Current file); Jan 16, 1991; ... the Justice Department\textquotesingle{}}
\end{Highlighting}
\end{Shaded}

A maior parte do texto está legível, mas há sequências estranhas, como
\texttt{} bem no início e vários \texttt{\textbackslash{}n} ao longo do
texto:

\begin{itemize}
\tightlist
\item
  A sequência \texttt{} é apenas uma declaração explícita (uma
  ``assinatura'') de que o arquivo foi codificado em UTF‑8. Nem todo
  arquivo UTF‑8 contém essa sequência.
\item
  As sequências \texttt{\textbackslash{}n} indicam \textbf{quebra de
  linha}.
\end{itemize}

Isso fica evidente se usarmos a função \texttt{print()} para imprimir os
primeiros 1000 caracteres:

\begin{Shaded}
\begin{Highlighting}[]
\CommentTok{\# Imprime os primeiros 1000 caracteres da variável \textquotesingle{}text\textquotesingle{}}
\BuiltInTok{print}\NormalTok{(text[:}\DecValTok{1000}\NormalTok{])}
\end{Highlighting}
\end{Shaded}

\begin{Shaded}
\begin{Highlighting}[]
\NormalTok{U.S. TAKING STEPS TO CURB TERRORISM: F.B.I. Is Ordered to Find Iraqis Whose Visas Have Expired}
\NormalTok{By JAMES BARRON}
\NormalTok{New York Times (1923{-}Current file); Jan 16, 1991;}
\NormalTok{ProQuest Historical Newspapers: The New York Times with Index pg. A15}
\NormalTok{...}
\NormalTok{   The Federal Bureau of Investigation has been ordered to track down as many as 3,000 Iraqis in this country whose visas have expired, the Justice Department said yesterday.}
\end{Highlighting}
\end{Shaded}

Como se vê, o Python sabe interpretar as sequências
\texttt{\textbackslash{}n} e insere uma quebra de linha ao encontrá‑las
durante a impressão. Note também que as primeiras linhas do arquivo
contêm \textbf{metadados} do artigo (nome, autor e fonte), que antecedem
o corpo do texto.

\begin{center}\rule{0.5\linewidth}{0.5pt}\end{center}

\subsection{4. Manipulando texto}\label{manipulando-texto}

Como todo o conteúdo guardado em \texttt{text} é um objeto
\emph{string}, podemos usar todos os métodos de manipulação de
\emph{strings} do Python.

\subsubsection{\texorpdfstring{4.1 Substituindo com
\texttt{replace()}}{4.1 Substituindo com replace()}}\label{substituindo-com-replace}

Vamos usar o método \texttt{replace()} para trocar todas as quebras de
linha \texttt{"\textbackslash{}n"} por \emph{strings} vazias \texttt{""}
e guardar o resultado em \texttt{processed\_text}:

\begin{Shaded}
\begin{Highlighting}[]
\CommentTok{\# Substitui as quebras de linha \textbackslash{}n por strings vazias e atribui o resultado}
\CommentTok{\# à variável \textquotesingle{}processed\_text\textquotesingle{}}
\NormalTok{processed\_text }\OperatorTok{=}\NormalTok{ text.replace(}\StringTok{\textquotesingle{}}\CharTok{\textbackslash{}n}\StringTok{\textquotesingle{}}\NormalTok{, }\StringTok{\textquotesingle{}\textquotesingle{}}\NormalTok{)}

\CommentTok{\# Imprime os primeiros 1000 caracteres da variável \textquotesingle{}processed\_text\textquotesingle{}}
\BuiltInTok{print}\NormalTok{(processed\_text[:}\DecValTok{1000}\NormalTok{])}
\end{Highlighting}
\end{Shaded}

Agora todo o texto ficou ``grudado''. Ainda assim, conseguimos
identificar o início de cada parágrafo, marcado por \textbf{três espaços
em branco}.

\begin{quote}
⚠ Atenção: remover as quebras de linha também fez os metadados do artigo
se fundirem em um único parágrafo. Sempre fique atento aos
\textbf{efeitos indesejados} de substituições e outras transformações!
\end{quote}

\subsubsection{\texorpdfstring{4.2 Dividindo com
\texttt{split()}}{4.2 Dividindo com split()}}\label{dividindo-com-split}

Se nos interessa apenas o corpo do artigo, podemos remover os metadados
com facilidade, pois sabemos que eles estão separados do corpo por três
espaços. A maneira mais simples é usar \texttt{split()} para
\textbf{dividir a \emph{string} em uma lista}, usando os três espaços
como separador:

\begin{Shaded}
\begin{Highlighting}[]
\CommentTok{\# Usa o método split() com três espaços como separador. Atribui o}
\CommentTok{\# resultado à variável \textquotesingle{}processed\_text\textquotesingle{}.}
\NormalTok{processed\_text }\OperatorTok{=}\NormalTok{ processed\_text.split(sep}\OperatorTok{=}\StringTok{\textquotesingle{}   \textquotesingle{}}\NormalTok{)}

\CommentTok{\# Imprime o resultado da variável \textquotesingle{}processed\_text\textquotesingle{}}
\BuiltInTok{print}\NormalTok{(processed\_text)}
\end{Highlighting}
\end{Shaded}

O método \texttt{split()} retorna uma \textbf{lista} de objetos
\emph{string}. Podemos confirmar isso verificando o tipo do objeto:

\begin{Shaded}
\begin{Highlighting}[]
\CommentTok{\# Verifica o tipo do objeto na variável \textquotesingle{}processed\_text\textquotesingle{}}
\BuiltInTok{type}\NormalTok{(processed\_text)}
\end{Highlighting}
\end{Shaded}

\begin{Shaded}
\begin{Highlighting}[]
\NormalTok{list}
\end{Highlighting}
\end{Shaded}

Os metadados ficaram no \textbf{primeiro item} da lista, pois a primeira
sequência de três espaços aparece justamente onde os metadados terminam.
Vamos buscar esse primeiro item --- lembre que o Python começa a contar
do zero, então ele está no índice \texttt{0}:

\begin{Shaded}
\begin{Highlighting}[]
\CommentTok{\# Recupera o objeto string no índice 0 da lista \textquotesingle{}processed\_text\textquotesingle{}}
\NormalTok{processed\_text[}\DecValTok{0}\NormalTok{]}
\end{Highlighting}
\end{Shaded}

\subsubsection{\texorpdfstring{4.3 Removendo itens com
\texttt{pop()}}{4.3 Removendo itens com pop()}}\label{removendo-itens-com-pop}

Para remover os metadados e manter apenas o corpo do texto, usamos o
método \texttt{pop()} da lista. Ele espera um número inteiro: o índice
do item a ser removido.

\begin{Shaded}
\begin{Highlighting}[]
\CommentTok{\# Chama o método pop() da lista \textquotesingle{}processed\_text\textquotesingle{} com o}
\CommentTok{\# índice do item a ser removido.}
\NormalTok{processed\_text.pop(}\DecValTok{0}\NormalTok{)}
\end{Highlighting}
\end{Shaded}

Você pode se perguntar por que \textbf{não} atribuímos o resultado a uma
variável. A resposta é que as listas do Python são \textbf{mutáveis}, ou
seja, podem ser alteradas ``no lugar'' (\emph{in place}). O método
\texttt{pop()} modifica a lista sem precisar reatribuir o valor à
variável. Podemos conferir recuperando os três primeiros itens:

\begin{Shaded}
\begin{Highlighting}[]
\CommentTok{\# Recupera os três primeiros itens da lista \textquotesingle{}processed\_text\textquotesingle{}}
\NormalTok{processed\_text[:}\DecValTok{3}\NormalTok{]}
\end{Highlighting}
\end{Shaded}

O primeiro item da lista \textbf{não corresponde mais} aos metadados.

\subsubsection{\texorpdfstring{4.4 Reunindo com
\texttt{join()}}{4.4 Reunindo com join()}}\label{reunindo-com-join}

Para converter a lista de volta em uma \emph{string}, usamos o método
\texttt{join()} de um objeto \emph{string}. O \texttt{join()} espera um
\textbf{iterável} como entrada (algo que possa ser percorrido, como uma
lista ou um dicionário).

Aqui pode haver certa confusão: o \texttt{join()} deve ser chamado
\textbf{sobre a \emph{string} que servirá de ``cola''} entre os itens do
iterável. No nosso caso, queremos usar como cola a sequência original
que separava os parágrafos --- uma quebra de linha e três espaços
(\texttt{\textquotesingle{}\textbackslash{}n\ \ \ \textquotesingle{}}):

\begin{Shaded}
\begin{Highlighting}[]
\CommentTok{\# Usa o método join() para unir os itens da lista \textquotesingle{}processed\_text\textquotesingle{}}
\CommentTok{\# usando a string \textquotesingle{}\textbackslash{}n   \textquotesingle{} — uma quebra de linha e três espaços. Guarda o}
\CommentTok{\# resultado na variável de mesmo nome.}
\NormalTok{processed\_text }\OperatorTok{=} \StringTok{\textquotesingle{}}\CharTok{\textbackslash{}n}\StringTok{   \textquotesingle{}}\NormalTok{.join(processed\_text)}

\CommentTok{\# Verifica o resultado imprimindo os primeiros 1000 caracteres da string}
\CommentTok{\# resultante em \textquotesingle{}processed\_text\textquotesingle{}}
\BuiltInTok{print}\NormalTok{(processed\_text[:}\DecValTok{1000}\NormalTok{])}
\end{Highlighting}
\end{Shaded}

Aplicar o \texttt{join()} devolve uma \emph{string} com as quebras de
parágrafo originais!

\subsubsection{\texorpdfstring{4.5 Um ``pipeline'' de substituições com
laço
\texttt{for}}{4.5 Um ``pipeline'' de substituições com laço for}}\label{um-pipeline-de-substituiuxe7uxf5es-com-lauxe7o-for}

Examinando o texto de perto, dá para ver resquícios da digitalização: o
OCR resultou em uma \textbf{mistura de aspas} de tipos diferentes, como
\texttt{"}, \texttt{“}, \texttt{”}, \texttt{’’} e \texttt{‘‘}. Se
quiséssemos extrair citações do corpo do texto, seria bom padronizar as
aspas. Vamos escolher \texttt{"} (uma aspa dupla simples) como padrão.

Poderíamos aplicar \texttt{replace()} separadamente para cada tipo de
aspa, mas isso seria tedioso. Para tornar o processo mais eficiente,
combinamos duas estruturas de dados do Python: \textbf{listas} e
\textbf{tuplas}.

Começamos definindo uma lista chamada \texttt{pipeline}. Criamos e
preenchemos uma lista simplesmente colocando objetos entre colchetes
\texttt{{[}{]}}, separados por vírgula. Como \texttt{replace()} recebe
duas \emph{strings}, combinamos cada par em uma \textbf{tupla} ---
estruturas finitas e ordenadas, marcadas por parênteses \texttt{(\ )}.
Em cada tupla, colocamos o caractere a ser substituído na primeira
\emph{string} e o substituto na segunda:

\begin{Shaded}
\begin{Highlighting}[]
\CommentTok{\# Define uma lista com quatro tuplas, cada uma com duas strings: o caractere}
\CommentTok{\# a ser substituido e o seu substituto.}
\NormalTok{pipeline }\OperatorTok{=}\NormalTok{ [(}\StringTok{\textquotesingle{}“\textquotesingle{}}\NormalTok{, }\StringTok{\textquotesingle{}"\textquotesingle{}}\NormalTok{), (}\StringTok{\textquotesingle{}´´\textquotesingle{}}\NormalTok{, }\StringTok{\textquotesingle{}"\textquotesingle{}}\NormalTok{), (}\StringTok{\textquotesingle{}”\textquotesingle{}}\NormalTok{, }\StringTok{\textquotesingle{}"\textquotesingle{}}\NormalTok{), (}\StringTok{\textquotesingle{}’’\textquotesingle{}}\NormalTok{, }\StringTok{\textquotesingle{}"\textquotesingle{}}\NormalTok{)]}
\end{Highlighting}
\end{Shaded}

Isso ilustra como diferentes estruturas de dados costumam ser
\textbf{aninhadas} no Python: a lista contém tuplas, e as tuplas contêm
\emph{strings}.

Agora podemos percorrer cada item da lista com um laço \texttt{for}, que
itera pelos itens na ordem em que aparecem. Cada item é uma tupla com
duas \emph{strings}. Para entrar no laço, o Python espera que a próxima
linha esteja \textbf{indentada} (use a tecla \texttt{Tab\ ↹}):

\begin{Shaded}
\begin{Highlighting}[]
\CommentTok{\# Percorre as tuplas da lista \textquotesingle{}pipeline\textquotesingle{}. Cada tupla tem dois valores, que}
\CommentTok{\# atribuimos as variaveis \textquotesingle{}old\textquotesingle{} e \textquotesingle{}new\textquotesingle{} automaticamente!}
\ControlFlowTok{for}\NormalTok{ old, new }\KeywordTok{in}\NormalTok{ pipeline:}

    \CommentTok{\# Usa o metodo replace() para substituir a string da variavel \textquotesingle{}old\textquotesingle{}}
    \CommentTok{\# pela string da variavel \textquotesingle{}new\textquotesingle{}}
\NormalTok{    processed\_text }\OperatorTok{=}\NormalTok{ processed\_text.replace(old, new)}
\end{Highlighting}
\end{Shaded}

O que acontece dentro do laço é exatamente o que fizemos antes com
\texttt{replace()}, mas, em vez de definir manualmente as
\emph{strings}, usamos as \emph{strings} contidas nas variáveis
\texttt{old} e \texttt{new}! A cada volta, atualizamos automaticamente a
\emph{string} em \texttt{processed\_text}.

\begin{Shaded}
\begin{Highlighting}[]
\CommentTok{\# Imprime a string}
\BuiltInTok{print}\NormalTok{(processed\_text)}
\end{Highlighting}
\end{Shaded}

Conseguimos realizar uma série de substituições percorrendo a lista de
tuplas, que definia os padrões a substituir e seus substitutos.

\textbf{Recapitulando a sintaxe do \texttt{for}:} declaramos o início do
laço com \texttt{for}, seguido de uma variável usada para se referir aos
itens obtidos da lista. A lista percorrida vem precedida de \texttt{in}
e do nome da variável atribuída à lista inteira.

Para entender melhor, vejamos o que acontece se definirmos
\textbf{apenas uma} variável, \texttt{our\_tuple}, para os itens
obtidos:

\begin{Shaded}
\begin{Highlighting}[]
\CommentTok{\# Percorre os itens da variavel \textquotesingle{}pipeline\textquotesingle{}}
\ControlFlowTok{for}\NormalTok{ our\_tuple }\KeywordTok{in}\NormalTok{ pipeline:}

    \CommentTok{\# Imprime o objeto retornado}
    \BuiltInTok{print}\NormalTok{(our\_tuple)}
\end{Highlighting}
\end{Shaded}

\begin{Shaded}
\begin{Highlighting}[]
\NormalTok{(\textquotesingle{}“\textquotesingle{}, \textquotesingle{}"\textquotesingle{})}
\NormalTok{(\textquotesingle{}´´\textquotesingle{}, \textquotesingle{}"\textquotesingle{})}
\NormalTok{(\textquotesingle{}”\textquotesingle{}, \textquotesingle{}"\textquotesingle{})}
\NormalTok{(\textquotesingle{}’’\textquotesingle{}, \textquotesingle{}"\textquotesingle{})}
\end{Highlighting}
\end{Shaded}

Isso imprime as \textbf{tuplas} inteiras! O Python é esperto o bastante
para entender que uma única variável se refere aos itens (as tuplas) da
lista, ao passo que, com duas variáveis, ele avança até as
\emph{strings} contidas dentro de cada tupla. Ao escrever laços
\texttt{for}, preste muita atenção aos itens contidos na lista!

\begin{center}\rule{0.5\linewidth}{0.5pt}\end{center}

\subsection{5. Manipulando texto em
escala}\label{manipulando-texto-em-escala}

O ideal é que o Python permita manipular texto \textbf{em escala}, isto
é, aplicar o mesmo procedimento a dez, cem ou mil arquivos com o mesmo
esforço. Para isso, precisamos definir padrões \textbf{mais flexíveis}
do que as \emph{strings} fixas usadas com \texttt{replace()}, além de
abrir e fechar arquivos automaticamente. Essas capacidades vêm dos
módulos de \textbf{expressões regulares} e de \textbf{manipulação de
arquivos}.

\subsubsection{\texorpdfstring{5.1 Expressões regulares (\emph{regular
expressions})}{5.1 Expressões regulares (regular expressions)}}\label{expressuxf5es-regulares-regular-expressions}

As expressões regulares são uma ``linguagem'' que permite definir
\textbf{padrões de busca}. Esses padrões podem ser usados para encontrar
(ou encontrar e substituir) trechos em objetos \emph{string}.
Diferentemente das \emph{strings} fixas, elas permitem definir
\textbf{curingas} (caracteres que representam qualquer caractere),
\textbf{quantificadores} (que casam sequências repetidas) e muito mais.

O Python oferece expressões regulares por meio do módulo \texttt{re},
que ativamos com o comando \texttt{import}:

\begin{Shaded}
\begin{Highlighting}[]
\ImportTok{import}\NormalTok{ re}
\end{Highlighting}
\end{Shaded}

Vamos carregar um arquivo, ler seu conteúdo, atribuir os
\textbf{últimos} 2000 caracteres à variável \texttt{extract} e imprimir
o resultado:

\begin{Shaded}
\begin{Highlighting}[]
\CommentTok{\# Define o caminho do arquivo e o abre para leitura (r) com codificacao utf{-}8}
\ControlFlowTok{with} \BuiltInTok{open}\NormalTok{(}\BuiltInTok{file}\OperatorTok{=}\StringTok{\textquotesingle{}data/WP\_1990{-}08{-}10{-}25A.txt\textquotesingle{}}\NormalTok{, mode}\OperatorTok{=}\StringTok{\textquotesingle{}r\textquotesingle{}}\NormalTok{, encoding}\OperatorTok{=}\StringTok{\textquotesingle{}utf{-}8\textquotesingle{}}\NormalTok{) }\ImportTok{as} \BuiltInTok{file}\NormalTok{:}

    \CommentTok{\# Le o conteudo do arquivo com o metodo .read()}
\NormalTok{    text }\OperatorTok{=} \BuiltInTok{file}\NormalTok{.read()}

\CommentTok{\# Pega os *ultimos* 2000 caracteres — note o sinal de menos antes do numero}
\NormalTok{extract }\OperatorTok{=}\NormalTok{ text[}\OperatorTok{{-}}\DecValTok{2000}\NormalTok{:]}

\CommentTok{\# Imprime o resultado}
\BuiltInTok{print}\NormalTok{(extract)}
\end{Highlighting}
\end{Shaded}

O texto tem muitos erros de OCR, principalmente sequências como
\texttt{....} e \texttt{,,,,}.

Vamos \textbf{compilar} nossa primeira expressão regular, que busca
sequências de dois ou mais pontos finais, com a função
\texttt{compile()} do módulo \texttt{re}. Ela recebe uma \emph{string}
como entrada. Note o prefixo \texttt{r} antes da \emph{string}: ele diz
ao Python para guardar a \emph{string} em formato ``bruto''
(\emph{raw}), ou seja, exatamente como aparece.

\begin{Shaded}
\begin{Highlighting}[]
\CommentTok{\# Compila uma expressao regular e a atribui a variavel \textquotesingle{}stops\textquotesingle{}}
\NormalTok{stops }\OperatorTok{=}\NormalTok{ re.}\BuiltInTok{compile}\NormalTok{(}\VerbatimStringTok{r\textquotesingle{}}\CharTok{\textbackslash{}.}\OperatorTok{\{2,\}}\VerbatimStringTok{\textquotesingle{}}\NormalTok{)}

\CommentTok{\# Vamos verificar o tipo da expressao regular!}
\BuiltInTok{type}\NormalTok{(stops)}
\end{Highlighting}
\end{Shaded}

\begin{Shaded}
\begin{Highlighting}[]
\NormalTok{re.Pattern}
\end{Highlighting}
\end{Shaded}

Vamos destrinchar essa expressão:

\begin{itemize}
\tightlist
\item
  Ela é definida por uma \emph{string} do Python (aspas simples
  \texttt{\textquotesingle{}\ \ \textquotesingle{}}).
\item
  Precisamos de uma \textbf{barra invertida} \texttt{\textbackslash{}}
  antes do ponto final \texttt{.}, porque, nas expressões regulares, o
  ponto final é um \textbf{curinga} que representa qualquer caractere. A
  barra diz ao Python que queremos o ponto final ``de verdade''.
\item
  As \textbf{chaves} \texttt{\{\ \}} instruem a expressão a buscar
  ocorrências do item anterior (\texttt{\textbackslash{}.}, nosso ponto
  de verdade) que aconteçam \textbf{duas ou mais vezes} (\texttt{2,}).
  Isso preserva os usos legítimos de um único ponto final.
\end{itemize}

Em português claro: buscamos ocorrências de \textbf{dois ou mais pontos
finais}.

Para aplicar a expressão a algum texto, usamos o método \texttt{sub()}
do objeto recém‑criado \texttt{stops}. O \texttt{sub()} recebe dois
argumentos:

\begin{itemize}
\tightlist
\item
  \texttt{repl}: a \emph{string} usada para \textbf{substituir} as
  ocorrências encontradas;
\item
  \texttt{string}: o objeto \emph{string} onde procurar as ocorrências.
\end{itemize}

O método retorna a \emph{string} modificada:

\begin{Shaded}
\begin{Highlighting}[]
\CommentTok{\# Aplica a expressao regular ao texto em \textquotesingle{}extract\textquotesingle{} e salva a saida}
\CommentTok{\# na mesma variavel, essencialmente sobrescrevendo o texto antigo.}
\NormalTok{extract }\OperatorTok{=}\NormalTok{ stops.sub(repl}\OperatorTok{=}\StringTok{\textquotesingle{}\textquotesingle{}}\NormalTok{, string}\OperatorTok{=}\NormalTok{extract)}

\CommentTok{\# Imprime o texto para examinar o resultado}
\BuiltInTok{print}\NormalTok{(extract)}
\end{Highlighting}
\end{Shaded}

As sequências de pontos finais desapareceram.

\subsubsection{5.2 Alternativas na expressão
regular}\label{alternativas-na-expressuxe3o-regular}

Podemos tornar a expressão mais poderosa acrescentando
\textbf{alternativas}. Vamos compilar outra e guardá‑la em
\texttt{punct}:

\begin{Shaded}
\begin{Highlighting}[]
\CommentTok{\# Compila uma expressao regular e a atribui a variavel \textquotesingle{}punct\textquotesingle{}}
\NormalTok{punct }\OperatorTok{=}\NormalTok{ re.}\BuiltInTok{compile}\NormalTok{(}\VerbatimStringTok{r\textquotesingle{}}\KeywordTok{(}\CharTok{\textbackslash{}.}\ControlFlowTok{|}\VerbatimStringTok{,}\KeywordTok{)}\OperatorTok{\{2,\}}\VerbatimStringTok{\textquotesingle{}}\NormalTok{)}
\end{Highlighting}
\end{Shaded}

A novidade são os \textbf{parênteses} \texttt{(\ )} e a \textbf{barra
vertical} \texttt{\textbar{}} entre eles, que separa o ponto final
\texttt{\textbackslash{}.} da vírgula \texttt{,}. Caracteres entre
parênteses e separados por \texttt{\textbar{}} marcam
\textbf{alternativas}. Em português claro: buscamos ocorrências de
\textbf{dois ou mais pontos finais ou vírgulas}.

Para garantir que o padrão funcione como esperado, recuperamos o texto
original de \texttt{text} e o reatribuímos a \texttt{extract},
sobrescrevendo as edições anteriores:

\begin{Shaded}
\begin{Highlighting}[]
\CommentTok{\# "Reinicia" a variavel extract pegando os ultimos 2000 caracteres da string original}
\NormalTok{extract }\OperatorTok{=}\NormalTok{ text[}\OperatorTok{{-}}\DecValTok{2000}\NormalTok{:]}

\CommentTok{\# Aplica a expressao regular}
\NormalTok{extract }\OperatorTok{=}\NormalTok{ punct.sub(repl}\OperatorTok{=}\StringTok{\textquotesingle{}\textquotesingle{}}\NormalTok{, string}\OperatorTok{=}\NormalTok{extract)}

\CommentTok{\# Imprime o resultado}
\BuiltInTok{print}\NormalTok{(extract)}
\end{Highlighting}
\end{Shaded}

Sucesso! As sequências de pontos finais \textbf{e} de vírgulas podem ser
removidas com uma única expressão regular.

Sequências mais irregulares vindas de erros de OCR --- como
\texttt{\textquotesingle{}-\textquotesingle{}*},
\texttt{-\textgreater{}."}, \texttt{/*—.} --- são bem mais difíceis de
capturar e exigiriam expressões mais complexas. Essa complexidade é
justamente o que torna as expressões regulares tão poderosas, mas
aprender a usá‑las leva \textbf{tempo e paciência}.

\begin{tipbox}

\textbf{Dica:} use um serviço como
\href{https://regex101.com}{regex101.com} para aprender e testar
expressões regulares interativamente.

\end{tipbox}

Na prática, criar expressões que cubram o máximo de casos é
particularmente difícil. Capturar a maioria dos erros --- talvez
distribuindo as transformações em uma \textbf{sequência de etapas} (um
\emph{pipeline}) --- já ajuda muito a preparar o texto para análise.
Lembre‑se, porém: para identificar padrões de manipulação de forma
confiável, você deve sempre observar \textbf{mais de um} texto do seu
corpus.

\subsubsection{5.3 Processando múltiplos
arquivos}\label{processando-muxfaltiplos-arquivos}

Muitos corpora contêm textos em vários arquivos. Para manipular grandes
volumes de texto com eficiência, precisamos abrir os arquivos, ler seu
conteúdo, executar as operações desejadas e fechá‑los --- tudo
\textbf{programaticamente}.

Esse procedimento fica bem simples com a classe \texttt{Path} do módulo
\texttt{pathlib} do Python. Usar \texttt{from\ ...\ import\ ...} permite
importar \textbf{apenas parte} de um módulo --- aqui, só a classe
\texttt{Path}:

\begin{Shaded}
\begin{Highlighting}[]
\ImportTok{from}\NormalTok{ pathlib }\ImportTok{import}\NormalTok{ Path}
\end{Highlighting}
\end{Shaded}

A classe \texttt{Path} codifica informações sobre caminhos em uma
estrutura de diretórios. O grande trunfo dela é \textbf{inferir
automaticamente} o tipo de caminho usado pelo seu sistema operacional.
Isso importa porque Windows, Linux e macOS usam caminhos de arquivo
diferentes; a classe \texttt{Path} nos poupa de muita dor de cabeça,
sobretudo se quisermos que o código rode em sistemas diferentes.

Nosso repositório contém um diretório chamado \texttt{data}, com os
arquivos de texto que viemos usando. Vamos inicializar um objeto
\texttt{Path} apontando para esse diretório e atribuí‑lo à variável
\texttt{corpus\_dir}:

\begin{Shaded}
\begin{Highlighting}[]
\CommentTok{\# Cria um objeto Path que aponta para o diretorio \textquotesingle{}data\textquotesingle{} e o atribui}
\CommentTok{\# a variavel \textquotesingle{}corpus\_dir\textquotesingle{}}
\NormalTok{corpus\_dir }\OperatorTok{=}\NormalTok{ Path(}\StringTok{\textquotesingle{}data\textquotesingle{}}\NormalTok{)}
\end{Highlighting}
\end{Shaded}

O objeto \texttt{Path} tem vários métodos e atributos úteis. Podemos,
por exemplo, verificar se o caminho é válido com \texttt{exists()}, que
retorna um valor \textbf{booleano} (\texttt{True} ou \texttt{False}):

\begin{Shaded}
\begin{Highlighting}[]
\CommentTok{\# Usa o metodo exists() para verificar se o caminho e valido}
\NormalTok{corpus\_dir.exists()}
\end{Highlighting}
\end{Shaded}

\begin{Shaded}
\begin{Highlighting}[]
\NormalTok{True}
\end{Highlighting}
\end{Shaded}

Também podemos checar se o caminho é um diretório com
\texttt{is\_dir()}, e garantir que ele \textbf{não} aponta para um
arquivo com \texttt{is\_file()}:

\begin{Shaded}
\begin{Highlighting}[]
\CommentTok{\# Usa o metodo is\_dir() para verificar se o caminho aponta para um diretorio}
\NormalTok{corpus\_dir.is\_dir()}
\end{Highlighting}
\end{Shaded}

\begin{Shaded}
\begin{Highlighting}[]
\NormalTok{True}
\end{Highlighting}
\end{Shaded}

\begin{Shaded}
\begin{Highlighting}[]
\CommentTok{\# Usa o metodo is\_file() para verificar se o caminho aponta para um arquivo}
\NormalTok{corpus\_dir.is\_file()}
\end{Highlighting}
\end{Shaded}

\begin{Shaded}
\begin{Highlighting}[]
\NormalTok{False}
\end{Highlighting}
\end{Shaded}

Sabendo que o caminho aponta para um diretório, usamos o método
\texttt{glob()} para coletar todos os arquivos de texto. O nome
\texttt{glob} vem de \emph{global} e foi implementado originalmente como
um programa para casar nomes de arquivos e caminhos usando curingas. O
\texttt{glob()} exige o argumento \texttt{pattern}, que recebe uma
\emph{string}: o \texttt{*} (asterisco) atua como \textbf{curinga},
representando qualquer sequência de caracteres antes de \texttt{.txt}
(sufixo comum de arquivos de texto simples). Convertendo o resultado em
lista com \texttt{list()}, podemos percorrê‑lo facilmente:

\begin{Shaded}
\begin{Highlighting}[]
\CommentTok{\# Coleta todos os arquivos com sufixo .txt no diretorio \textquotesingle{}corpus\_dir\textquotesingle{} e converte o resultado em uma lista}
\NormalTok{files }\OperatorTok{=} \BuiltInTok{list}\NormalTok{(corpus\_dir.glob(pattern}\OperatorTok{=}\StringTok{\textquotesingle{}*.txt\textquotesingle{}}\NormalTok{))}

\CommentTok{\# Mostra o resultado}
\NormalTok{files}
\end{Highlighting}
\end{Shaded}

\begin{Shaded}
\begin{Highlighting}[]
\NormalTok{[PosixPath(\textquotesingle{}data/WP\_1990{-}08{-}10{-}25A.txt\textquotesingle{}),}
\NormalTok{ PosixPath(\textquotesingle{}data/NYT\_1991{-}01{-}16{-}A15.txt\textquotesingle{}),}
\NormalTok{ PosixPath(\textquotesingle{}data/WP\_1991{-}01{-}17{-}A1B.txt\textquotesingle{})]}
\end{Highlighting}
\end{Shaded}

Temos uma lista de três objetos \texttt{Path} apontando para três
arquivos! Isso nos permite percorrê‑los com um laço \texttt{for} e
manipular o texto de cada um. No bloco abaixo, iteramos sobre cada
arquivo, lemos e modificamos seu conteúdo e o gravamos em um novo
arquivo:

\begin{Shaded}
\begin{Highlighting}[]
\CommentTok{\# Percorre a lista de objetos Path em \textquotesingle{}files\textquotesingle{}. Refere{-}se a cada arquivo}
\CommentTok{\# pela variavel \textquotesingle{}file\textquotesingle{}.}
\ControlFlowTok{for} \BuiltInTok{file} \KeywordTok{in}\NormalTok{ files:}

    \CommentTok{\# Usa o metodo read\_text() de um objeto Path para ler o conteudo do arquivo.}
    \CommentTok{\# Passa o valor \textquotesingle{}utf{-}8\textquotesingle{} ao argumento \textquotesingle{}encoding\textquotesingle{} para declarar a codificacao.}
    \CommentTok{\# Guarda o resultado na variavel \textquotesingle{}text\textquotesingle{}.}
\NormalTok{    text }\OperatorTok{=} \BuiltInTok{file}\NormalTok{.read\_text(encoding}\OperatorTok{=}\StringTok{\textquotesingle{}utf{-}8\textquotesingle{}}\NormalTok{)}

    \CommentTok{\# Aplica a expressao regular definida acima para remover a pontuacao}
    \CommentTok{\# excessiva do texto. Guarda o resultado na variavel \textquotesingle{}mod\_text\textquotesingle{}.}
\NormalTok{    mod\_text }\OperatorTok{=}\NormalTok{ punct.sub(}\StringTok{\textquotesingle{}\textquotesingle{}}\NormalTok{, text)}

    \CommentTok{\# Define um novo nome de arquivo com o prefixo \textquotesingle{}mod\_\textquotesingle{} criando uma nova string.}
    \CommentTok{\# O objeto Path guarda o nome do arquivo como string no atributo \textquotesingle{}name\textquotesingle{}.}
    \CommentTok{\# Combina as duas strings com o operador \textquotesingle{}+\textquotesingle{}.}
\NormalTok{    new\_filename }\OperatorTok{=} \StringTok{\textquotesingle{}mod\_\textquotesingle{}} \OperatorTok{+} \BuiltInTok{file}\NormalTok{.name}

    \CommentTok{\# Define um novo objeto Path que aponta para o novo arquivo. O objeto Path}
    \CommentTok{\# junta automaticamente o diretorio e o nome do arquivo para nos.}
\NormalTok{    new\_path }\OperatorTok{=}\NormalTok{ Path(}\StringTok{\textquotesingle{}data\textquotesingle{}}\NormalTok{, new\_filename)}

    \CommentTok{\# Imprime uma mensagem de status usando formatacao de string. Ao adicionar}
    \CommentTok{\# o prefixo \textquotesingle{}f\textquotesingle{} a uma string, podemos usar chaves \{\} para inserir uma}
    \CommentTok{\# variavel dentro da string. Aqui inserimos o caminho atual do arquivo.}
    \BuiltInTok{print}\NormalTok{(}\SpecialStringTok{f\textquotesingle{}Writing modified text to }\SpecialCharTok{\{}\NormalTok{new\_path}\SpecialCharTok{\}}\SpecialStringTok{\textquotesingle{}}\NormalTok{)}

    \CommentTok{\# Usa o metodo write\_text() para escrever o texto modificado de \textquotesingle{}mod\_text\textquotesingle{}}
    \CommentTok{\# no arquivo, com codificacao UTF{-}8.}
\NormalTok{    new\_path.write\_text(mod\_text, encoding}\OperatorTok{=}\StringTok{\textquotesingle{}utf{-}8\textquotesingle{}}\NormalTok{)}
\end{Highlighting}
\end{Shaded}

\begin{Shaded}
\begin{Highlighting}[]
\NormalTok{Writing modified text to data/mod\_WP\_1990{-}08{-}10{-}25A.txt}
\NormalTok{Writing modified text to data/mod\_NYT\_1991{-}01{-}16{-}A15.txt}
\NormalTok{Writing modified text to data/mod\_WP\_1991{-}01{-}17{-}A1B.txt}
\end{Highlighting}
\end{Shaded}

Como se vê, os objetos \texttt{Path} oferecem dois métodos convenientes
para trabalhar com arquivos de texto: \texttt{read\_text()} e
\texttt{write\_text()}. Eles permitem ler e gravar texto \textbf{sem} a
instrução \texttt{with} --- e, assim como o \texttt{with}, o arquivo
apontado pelo \texttt{Path} é fechado automaticamente após a leitura.

\begin{center}\rule{0.5\linewidth}{0.5pt}\end{center}

\subsection{Quiz}\label{quiz}

Marque a alternativa correta (a resposta certa está destacada com ✅).

\textbf{1. Qual formato de texto é o mais interoperável para
programação?}

\begin{enumerate}
\def\labelenumi{\arabic{enumi}.}
\tightlist
\item
  Texto formatado (\emph{rich text})
\item
  Texto estruturado (\emph{structured text})
\item
  Texto simples (\emph{plain text}) ✅
\end{enumerate}

\textbf{2. Num texto, a sequência de dois caracteres ``barra‑n''
representa:}

\begin{enumerate}
\def\labelenumi{\arabic{enumi}.}
\tightlist
\item
  Um espaço em branco
\item
  Uma quebra de linha ✅
\item
  Uma tabulação
\end{enumerate}

\textbf{3. O UTF‑8 é retrocompatível com qual codificação?}

\begin{enumerate}
\def\labelenumi{\arabic{enumi}.}
\tightlist
\item
  ASCII ✅
\item
  Latin‑1
\item
  UTF‑16
\end{enumerate}

\textbf{4. Qual método de string divide o texto em uma lista?}

\begin{enumerate}
\def\labelenumi{\arabic{enumi}.}
\tightlist
\item
  \texttt{join()}
\item
  \texttt{split()} ✅
\item
  \texttt{replace()}
\end{enumerate}

\textbf{5. Numa expressão regular, o quantificador ``dois ou mais'' é
escrito como:}

\begin{enumerate}
\def\labelenumi{\arabic{enumi}.}
\tightlist
\item
  \texttt{\{2\}}
\item
  \texttt{\{2,\}} ✅
\item
  \texttt{\{,2\}}
\end{enumerate}

\textbf{6. Qual método da classe \texttt{Path} coleta arquivos por um
padrão com curinga (ex.: todos os \texttt{.txt})?}

\begin{enumerate}
\def\labelenumi{\arabic{enumi}.}
\tightlist
\item
  \texttt{glob()} ✅
\item
  \texttt{read\_text()}
\item
  \texttt{exists()}
\end{enumerate}

\subsection{Resumo da unidade}\label{resumo-da-unidade}

Nesta unidade, você aprendeu a:

\begin{enumerate}
\def\labelenumi{\arabic{enumi}.}
\tightlist
\item
  \textbf{distinguir} texto formatado, estruturado e simples, e a
  entender por que o texto simples é o formato preferido em programação;
\item
  \textbf{compreender} a codificação de caracteres (ASCII e
  Unicode/UTF‑8) e por que ela causa problemas;
\item
  \textbf{carregar} arquivos com \texttt{open()} + \texttt{with} +
  \texttt{read()}, lidando com marcas como \texttt{} e
  \texttt{\textbackslash{}n};
\item
  \textbf{manipular} \emph{strings} com \texttt{replace()},
  \texttt{split()}, \texttt{pop()} e \texttt{join()}, e montar um
  \emph{pipeline} de substituições com listas, tuplas e laços
  \texttt{for};
\item
  \textbf{definir padrões flexíveis} com expressões regulares
  (\texttt{re.compile()}, \texttt{sub()}, curingas, quantificadores e
  alternativas);
\item
  \textbf{processar vários arquivos} de uma vez com
  \texttt{pathlib.Path}, \texttt{glob()}, \texttt{read\_text()} e
  \texttt{write\_text()}.
\end{enumerate}

\subsubsection{Exercícios sugeridos}\label{exercuxedcios-sugeridos}

\begin{enumerate}
\def\labelenumi{\arabic{enumi}.}
\tightlist
\item
  Carregue um dos arquivos do diretório \texttt{data} e conte quantas
  quebras de linha (\texttt{\textbackslash{}n}) ele contém antes de
  qualquer limpeza.
\item
  Escreva um \emph{pipeline} (lista de tuplas) que padronize
  \textbf{travessões} (\texttt{—}, \texttt{–}, \texttt{-}) em um único
  caractere.
\item
  Crie uma expressão regular que remova sequências de dois ou mais
  espaços em branco, substituindo‑as por um único espaço.
\item
  Adapte o laço da Seção 5.3 para gravar os arquivos modificados em um
  novo diretório chamado \texttt{output} em vez de \texttt{data}.
\end{enumerate}

\begin{center}\rule{0.5\linewidth}{0.5pt}\end{center}

\subsubsection{Referências}\label{referuxeancias}

\begin{itemize}
\tightlist
\item
  Rögnvaldsson, E. et al.~\emph{Applied Language Technology} (MOOC).
  Universidade de Helsinque.
  \url{https://applied-language-technology.mooc.fi/}
\item
  Hill, M. J.; Hengchen, S. (2019). Quantifying the impact of dirty OCR
  on historical text analysis. \emph{Digital Scholarship in the
  Humanities}.
\end{itemize}

\end{document}
